\section{总结与延伸}

\subsection{核心观点回顾}

\begin{enumerate}
    \item \textbf{显式或隐式的 Plan 是必要且有效的}：对于语言模型求解复杂的长周期问题，Plan 机制可以显著提升表现。
    \item \textbf{Top-Down 与 Bottom-Up 的结合}：Top-Down 做任务分解形成 Plan，Bottom-Up 通过真实交互的 feedback 进行 Re-Plan，弥合头脑规划与真实执行之间的 Gap。
    \item \textbf{感知模型的能力与特点，扬长避短}：设计合适的 Workflow 和 Tool 来激发模型能力，规避其短板。
    \item \textbf{系统级的设计}：语言模型 + Tools + Prompts 需要统筹考虑，形成稳定高效的解决方案。
\end{enumerate}

\subsection{延伸思考}

\begin{itemize}
    \item \textbf{Workflow 设计的重要性}：UCLA 的 IMO 方案证明，模型本身有能力（capable）解决问题，关键在于 workflow 的设计是否能最大化地激发模型潜力。
    \item \textbf{从 Agentic Workflow 到 Autonomous Agent}：目前大量 Coding Agent（Claude Code、Codex、Cursor）都采用预定义的 workflow 而非完全自主的 Agent，因为可控、可靠、可预测。完全自主的 Agent 是终极目标，但受限于当前模型能力。
    \item \textbf{Agent 的设计和开发可能与模型训练同等重要}：从实用角度来看，通过精心设计 Agent 来达到更好的 performance，可能比训练更强的模型更为经济高效。
\end{itemize}

%% --- Fin del contenido --- %%

\end{document}
